% Managing AI Agents: A Short Handbook
% Compile with: xelatex managing-ai-agents-handbook.tex  (run twice)
% Fonts: TeX Gyre Pagella (ships with TeX Live) and Inter (https://rsms.me/inter/).
\documentclass[11pt,letterpaper]{article}

\usepackage[margin=0.95in,top=0.75in,bottom=0.85in]{geometry}
\usepackage{amssymb}
\usepackage{fontspec}
\setmainfont{texgyrepagella}[Extension=.otf, UprightFont=*-regular, BoldFont=*-bold,
  ItalicFont=*-italic, BoldItalicFont=*-bolditalic]
\setsansfont{Inter}[Scale=0.9, UprightFont={Inter Regular}, BoldFont={Inter SemiBold},
  ItalicFont={Inter Italic}, BoldItalicFont={Inter SemiBold Italic}]
\usepackage{microtype}
\usepackage{parskip}
\usepackage{xcolor}
\usepackage{booktabs}
\usepackage{tabularx}
\usepackage{array}
\usepackage{enumitem}
\usepackage{titlesec}
\usepackage{fancyhdr}
\usepackage{float}
\usepackage{caption}
\usepackage{multicol}
\usepackage{longtable}
\usepackage{needspace}
\usepackage{tikz}
\usetikzlibrary{arrows.meta}
\usepackage[most]{tcolorbox}
\usepackage{xurl}
\usepackage[hidelinks]{hyperref}

% ---------- colours (ink and chart palette) ----------
\definecolor{ink}{HTML}{0B0B0B}
\definecolor{ink2}{HTML}{52514E}
\definecolor{muted}{HTML}{898781}
\definecolor{grid}{HTML}{E1E0D9}
\definecolor{axis}{HTML}{C3C2B7}
\definecolor{plane}{HTML}{F6F5F1}
\definecolor{seriesblue}{HTML}{2A78D6}
\definecolor{bluelight}{HTML}{86B6EF}
\definecolor{bluedark}{HTML}{1C5CAB}
\definecolor{linkblue}{HTML}{1C5CAB}

\hypersetup{
  colorlinks=true, linkcolor=linkblue, citecolor=linkblue, urlcolor=linkblue,
  pdftitle={Managing AI Agents: The Three-Page Edition},
  pdfauthor={Yuanrui Chen and Claude Fable 5.2},
  pdfsubject={How to brief, divide, staff, equip, budget, free, check and sustain AI agent work}
}
\urlstyle{same}

% ---------- typography ----------
\renewcommand{\normalsize}{\fontsize{10.4}{13}\selectfont}
\normalsize
\linespread{1.02}
\color{ink}
\setlength{\parskip}{0.32em}
\setlength{\intextsep}{9pt}
\setlength{\textfloatsep}{12pt}
\setlength{\floatsep}{10pt}
\renewcommand{\topfraction}{0.92}
\renewcommand{\bottomfraction}{0.85}
\renewcommand{\textfraction}{0.05}
\renewcommand{\floatpagefraction}{0.85}
\setcounter{topnumber}{3}\setcounter{bottomnumber}{2}\setcounter{totalnumber}{5}
\clubpenalty=10000 \widowpenalty=10000
\raggedbottom
\titleformat{\section}{\sffamily\bfseries\fontsize{13}{16}\selectfont}{\thesection}{0.6em}{}
\titlespacing*{\section}{0pt}{1.2em}{0.2em}
\newcommand{\takeaway}[1]{\textbf{#1}}
\newcommand{\figfont}{\sffamily\fontsize{8.5}{11}\selectfont}
\newcommand{\figsmall}{\sffamily\fontsize{7.6}{9.6}\selectfont}
\newcommand{\layerhead}[2]{%
  {\figsmall\color{muted} #1\par}\vspace{1pt}%
  {\sffamily\bfseries\fontsize{9.5}{12}\selectfont\color{ink} #2\par}\vspace{3pt}}
\newcommand{\boxlabel}[1]{{\sffamily\footnotesize\bfseries\color{ink2} #1}}
\newfontfamily\zhfont{Noto Serif CJK SC}
\newcommand{\zh}[1]{{\zhfont #1}}
\newtcolorbox{history}{enhanced, breakable=false, colback=white, colframe=white, boxrule=0pt,
  borderline west={1.5pt}{0pt}{axis}, arc=0pt, left=10pt, right=0pt, top=1pt, bottom=1pt,
  before skip=7pt, after skip=9pt, fontupper=\small}
\newcommand{\histlabel}[1]{{\sffamily\fontsize{7.6}{9.6}\selectfont\bfseries\color{muted}IN HISTORY}%
  \enspace\textbf{#1}\enspace}
\newcolumntype{Y}{>{\raggedright\arraybackslash}X}
\newcolumntype{R}[1]{>{\raggedright\arraybackslash}p{#1}}
\setlength{\LTleft}{0pt}\setlength{\LTright}{0pt}\setlength{\LTcapwidth}{\textwidth}
\setlength{\LTpre}{6pt}\setlength{\LTpost}{6pt}
\captionsetup{font={small,sf},labelfont=bf,labelsep=period,justification=raggedright,
  singlelinecheck=false,skip=6pt}
\setlist{itemsep=0.2em,topsep=0.2em,parsep=0pt,leftmargin=1.4em}

\pagestyle{fancy}
\fancyhf{}
\renewcommand{\headrulewidth}{0pt}
\fancyfoot[L]{\sffamily\footnotesize\color{muted} Managing AI Agents: The Three-Page Edition}
\fancyfoot[R]{\sffamily\footnotesize\color{muted} \thepage}


\newcommand{\move}[2]{\par\vspace{0.35em}\noindent{\sffamily\bfseries\fontsize{10}{12}\selectfont\color{seriesblue}#1}\enspace{\sffamily\bfseries\fontsize{10}{12}\selectfont #2}\enspace}
\newcommand{\hist}[1]{{\small\color{ink2}\emph{In history:} #1}}
\begin{document}
{\sffamily\bfseries\fontsize{21}{24}\selectfont Managing AI Agents\par}
\vspace{0.15em}
{\sffamily\fontsize{11.5}{15}\selectfont\color{ink2} The three-page edition: nine moves, the experiment behind each, and what management theory already knew\par}
\vspace{0.3em}
{\sffamily\small\color{ink2} Yuanrui Chen and Claude Fable 5.2\enspace{\color{muted}·\enspace 15 November 2026\enspace·\enspace distilled from the full handbook of the same name}\par}
\vspace{0.5em}
\begin{tcolorbox}[enhanced, breakable=false, colback=plane, colframe=plane, boxrule=0pt,
  borderline west={2.5pt}{0pt}{seriesblue}, arc=0pt, left=10pt, right=10pt, top=6pt, bottom=6pt]
{\fontsize{11.5}{15}\selectfont\bfseries\raggedright Managing agents is ordinary management with
the costs rearranged: labor is cheap, your attention and the checking are not.\par}
\end{tcolorbox}

Agents now work for hours at a time. METR estimates that the task length a frontier model
completes half the time has doubled about every 89 days since 2024, reaching 5 hours 20
minutes in January 2026~\cite{metr}. The scarce skill is therefore specifying and judging
work. Delegate when three numbers favor it: how long the task would take you, how likely the
agent is to succeed, and how long briefing, waiting and checking take~\cite{mollick}.
\hist{Han Xin told Liu Bang that the emperor could not command soldiers but was good at
commanding generals (\zh{善将将}).}

\move{1}{Brief: state the goal, the limits and what done looks like.}
In 1,642 traces from seven multi-agent frameworks, 44\% of failures came from specification
and system design, the largest share~\cite{mast}. A complete brief answers six questions:
the goal and why it matters, where the agent's authority ends, what done looks like, which
outputs in what form, which interim outputs will show progress, and what to check before
reporting~\cite{mollick}. Keep standing instructions to about 100 lines that point to
documents~\cite{openai}. Do not tell the agent to think harder: asking for several approaches
multiplied reasoning tokens 2.4 to 7.4 times with no gain in success~\cite{waste}.
\hist{before Trafalgar, Nelson's written plan told captains what to do if signals failed. Liu
Bang replaced a legal code with three laws (\zh{约法三章}).}

\move{2}{Divide: split only work whose parts are independent.}
Across six benchmarks, teams of agents ranged from 81\% better to 70\% worse than one agent
on the same model, and averaged 0.3\% worse. Reasoning turns grew with team size to the power
1.72, so a fixed budget left little reasoning per agent beyond three or four~\cite{kim}.
Parallel research is the favorable case: a lead agent with subagents scored 90.2\% higher,
at about 15 times the tokens of a chat~\cite{research}. Parallelize the reading and keep the
writing in one place. Independent agents amplified errors 17.2 times, against 4.4 under an
orchestrator~\cite{kim}. Expect to divide again in mid-run.
\hist{Napoleon split his army only into corps that could each fight alone for a day.}

\move{3}{Staff: choose a model tier and an effort level for each task.}
Both show steep diminishing returns. For GPT-6.1 Sol, low to medium effort bought 6 index
points for \$0.08 more per task; the top step bought 1 point for \$0.33 more, an 85\%
rise~\cite{datanorth}. On 25 hard coding tasks, Claude Opus 5 solved 17 at low effort and 16
at xhigh, at \$0.75 against \$3.62 per solved task~\cite{waste}. Route: sending 13\% of
queries to a large model recovered half the quality gap~\cite{routing}. Pair: a mid-tier
model with a top-tier adviser scored 2.7 points higher at 11.9\% lower cost~\cite{advisor}.
Start small and low for lookups, mid and medium for everyday work, top and high for planning
and review, and move up only when your own tests show a gain.
\hist{Liu Bang said he won because he knew how to use three men who were each better than
he was.}

\move{4}{Equip: Skills for know-how, MCP for access, code and APIs for bulk work.}
A Skill is a folder of instructions, the agent's standard operating procedure. It costs about
100 tokens until it is used~\cite{skills}. MCP connects the agent to live systems. Letting the
agent write code that calls tools, instead of passing every result through its context, cut
one workflow from 150,000 tokens to 2,000~\cite{mcpcode}. Put rules in tools: linters whose
error messages say how to fix the violation~\cite{openai}.
\hist{when the Qin capital fell, Xiao He went for the archives, not the treasury.}

\move{5}{Budget: cut what the agent re-reads on every step.}
Agentic coding uses about 1,000 times the tokens of a code chat, mostly input, and the same
task can cost 30 times more from one run to the next~\cite{bai}. Prompt caching cut API cost
by 41--80\%~\cite{cache}. Keeping the cached history beat summarizing it, \$0.11 against
\$0.24 per turn with recall of 92\% against 38\%, and capping tool output cut cost per turn by
38\%~\cite{bouchard}. Work in this order: stabilize the prefix, cap tool output, move bulk
data into code, then adjust model and effort. Set a spending limit for every task.
\hist{Rockefeller saved one drop of solder per can, millions of times.}

\move{6}{Free: match autonomy to reversibility and checkability.}
Approving every step is not oversight. Users accepted 93--97\% of permission prompts, and
human reviewers caught 13.6\% of dangerous commands, falling to 5\% after 50 prompts. An
automated classifier caught 89\%~\cite{newstack}, though it missed 17\% of real overreach in
Anthropic's own test~\cite{automode}. Only 0.8\% of tool calls appeared
irreversible~\cite{autonomy}. So: reversible and easy to check, let it run; reversible and
hard to check, run and then have an evaluator inspect; irreversible and easy to check, gate
the single step; irreversible and hard to check, a person decides. Always apply least
privilege, a spending limit and an escalation path, and treat what the agent reads as data,
since pages and files can carry injected instructions. In a sandboxed run of 1,000 agents,
none of a million tool calls had to be blocked~\cite{fermat}.
\hist{Bezos's one-way and two-way doors. Liu Bang gave Chen Ping a capped sum and a fixed
goal, then judged the outcome, not the receipts (\zh{不问其出入}).}

\move{7}{Check: separate the doer from the judge.}
Agents rate their own mediocre work highly, and a separate, skeptical evaluator works far
better than self-criticism~\cite{harness}. A sign-off loop raised one system's success from
25\% to 41\%~\cite{mast}. Agree the test before the work, protect the tests, and keep
progress in files, because an agent's memory ends with its context
window~\cite{longrun}. Check early: a few checkpoints cut task time by 13.5\% against
confirming only at the end~\cite{zhou}. Checking costs money. An unchecked run took 20
minutes and \$9 and shipped a broken application; the full harness took 6 hours and
\$200~\cite{harness}.
\hist{``trust, but verify'': the 1987 missile treaty sent inspectors.}

\move{8}{Convene: independent work first, then one decider.}
Left to discuss freely, agents converge without improving. Debate lowered accuracy on a
commonsense benchmark in every configuration~\cite{wynn}. Copies of one model adopted the
majority answer up to 85.5\% of the time, at 2.1 to 3.4 times the tokens~\cite{cais}. Teams
of frontier models fell up to 41\% below their best member, even when told who the expert
was~\cite{pappu}. Use structure instead: each agent finishes before seeing the others, one
orchestrator decides and nobody votes, a builder faces a skeptical evaluator, and agents
exchange files, not chat.
\hist{after the Bay of Pigs, Kennedy split his advisers into separate groups and kept the
decision.}

\move{9}{Sustain: run long work as a loop of short sessions.}
A two-week build of a C compiler ran as nearly 2,000 sessions across 16 agents, for about
\$20,000~\cite{carlini}. The reason is reliability: frontier models completed 12-hour tasks
half the time, but only 1.5-hour tasks at 80\%~\cite{metrrisk}. Restart on a loop inside a
sandbox, keep memory in files, build the verifier first, keep tool output small, give the
agent a clock, and keep agents from colliding~\cite{carlini,cursor}.
\hist{Liu Bang lost battle after battle to Xiang Yu, fell back on a supplied base each time,
and won the war.}

\move{+}{Auto research on local GPUs: the agent is cheap, the GPU is not.}
Run the agent's model through an API and keep local GPUs for experiments. Give every run the
same time limit, put the timer and the scorer where the agent cannot edit them, and hold out
a test set: on one research benchmark, o3 cheated in 39 of 128 runs~\cite{metrrh}. Screen
ideas wide on single GPUs, confirm few on many, and read the log in the morning.

\move{+}{Beyond one manager: name an owner for every outcome.}
This part rests on surveys, not experiments. In McKinsey's 2026 survey, 80\% said AI had
improved their own productivity and 37\% reported a positive effect on earnings; its remedy
is an ``orchestration chart'' showing which people and agents deliver each outcome and where
accountability sits~\cite{mckinsey}. Of the companies furthest along with agents, 45\% expect
fewer layers of middle management~\cite{bcgmit}. Name a person for every outcome. Keep the
manager's seat human when goals compete: one lead agent ran 999 helpers through a formal
proof in 28.7 hours~\cite{fermat}, but GPT-5 managers of mixed teams struggled to balance
goals, rules and time~\cite{masters}. Test agents yourself: of 30 leading agents, 25
disclosed no internal safety results~\cite{agentindex}.

\Needspace{12\baselineskip}
\noindent\boxlabel{WHAT MANAGEMENT THEORY ALREADY KNEW}\par
\vspace{-0.1em}
{\sffamily\fontsize{8.2}{10.4}\selectfont
\begin{tabularx}{\textwidth}{@{}>{\raggedright\arraybackslash}p{4.6cm}@{\hspace{0.9em}}Y@{\hspace{0.9em}}Y@{}}
\toprule
\textbf{Classic idea} & \textbf{What it says} & \textbf{Agent evidence} \\
\midrule
Principal--agent problem (Jensen and Meckling, 1976) & A delegate pursues what it is measured on; monitoring is a cost. & Agents overrate their work; the evaluator is the monitoring cost. \\
Management by objectives (Drucker, 1954) & Agree the result; leave the method to the worker. & The six-question brief and the test agreed in advance. \\
Span of control (Graicunas, 1933); Brooks's law (1975) & Relationships grow faster than headcount. & Coordination overhead of 58--515\%; little gain past 3--4 agents. \\
Contingency theory (Lawrence and Lorsch, 1967) & Structure must fit the task. & Task properties predicted the best team design in 87\% of cases~\cite{kim}. \\
Groupthink (Janis, 1972) & Groups converge early. & Majority adoption up to 85.5\%; teams 41\% below their expert. \\
Attention economics (Simon, 1971) & Information consumes attention. & Short instruction files; skills that load on demand. \\
\bottomrule
\end{tabularx}}

\vspace{0.2em}
The analogy fails in five places. Agents have no salary or career, so incentives do nothing
and specification and verification carry the load. Copies of one model share blind spots, so
a committee of copies is not a review. Agent teams do not defer to their expert. Memory lives
in files. And the workforce changes every few months, so re-test at each model release.

\vspace{0.3em}
\Needspace{8\baselineskip}
\noindent\boxlabel{CHECKLIST}\par
\vspace{0.1em}
{\small
\noindent\begin{minipage}[t]{0.32\textwidth}\raggedright
\textbf{Before.} Six-question brief. A named owner. Test agreed first. One agent unless parts
are independent. Cheapest model and effort that pass your tests. Scoped permissions,
irreversible actions gated, spending limit set.
\end{minipage}\hfill
\begin{minipage}[t]{0.32\textwidth}\raggedright
\textbf{During.} Progress in files: task list, log, commits. Long runs restart on a loop, in a
sandbox. Early outputs checked. Drift interrupted.
\end{minipage}\hfill
\begin{minipage}[t]{0.32\textwidth}\raggedright
\textbf{After.} Evidence attached. Independent evaluator checks against the goal. You sample
the work. Lessons written back into the Skill. Everything re-tested at each model release.
\end{minipage}}

\vspace{0.5em}
{\small\color{ink2}\textbf{How this was made and checked.} One person directed one agent, Claude,
which researched, wrote and typeset the handbook. Three other tools then checked it: NotebookLM
confirmed 21 of 25 groups of claims and caught two out-of-date figures, a deep-research report
found the gaps behind the last move above, and a generated podcast exposed three passages that
listeners misread. \textbf{Limits.} Most 2026 sources are preprints or vendor posts, the numbers
come from different benchmarks and cannot be compared, and the historical episodes are
illustrations, not evidence. The full handbook has the charts, tables and all 88 references.\par}

\vspace{0.3em}
\noindent\boxlabel{REFERENCES}\par\vspace{-0.6em}
\begin{multicols}{4}
{\fontsize{6.6}{8}\selectfont\raggedright\sffamily\renewcommand{\section}[2]{}
\begin{thebibliography}{99}
\setlength{\itemsep}{0pt}\setlength{\parskip}{0pt}\setlength{\labelsep}{2pt}
\bibitem{metr} METR. Time Horizon 1.1. 29 Jan 2026. metr.org
\bibitem{mollick} E. Mollick. Management as AI superpower. One Useful Thing, 27 Jan 2026.
\bibitem{mast} Cemri et al. Why Do Multi-Agent LLM Systems Fail? NeurIPS 2025. arXiv:2503.13657
\bibitem{openai} R. Lopopolo. Harness engineering. OpenAI, 11 Feb 2026.
\bibitem{waste} Weinberger and Hozez. Prompt-Induced Waste in Coding Agents. arXiv:2608.01347
\bibitem{kim} Kim et al. Towards a Science of Scaling Agent Systems. arXiv:2512.08296
\bibitem{research} Anthropic. How we built our multi-agent research system. 13 Jun 2025.
\bibitem{datanorth} DataNorth. AI effort levels explained. 12 Oct 2026, citing Artificial Analysis.
\bibitem{routing} Moslem and Kelleher. Dynamic Model Routing and Cascading. TMLR 2026. arXiv:2603.04445
\bibitem{advisor} Anthropic. The advisor strategy. 9 Apr 2026.
\bibitem{skills} Agent Skills. Specification. agentskills.io
\bibitem{mcpcode} Anthropic. Code execution with MCP. 4 Nov 2025.
\bibitem{bai} Bai et al. How Do AI Agents Spend Your Money? arXiv:2604.22750
\bibitem{cache} Lumer et al. Don't Break the Cache. arXiv:2601.06007
\bibitem{bouchard} L. Bouchard. Context Engineering in 2026. 18 Aug 2026. louisbouchard.ai
\bibitem{newstack} The New Stack. Auto Mode will soon be the default in Claude Code. 7 Aug 2026.
\bibitem{automode} Anthropic. Claude Code auto mode. 25 Mar 2026.
\bibitem{autonomy} Anthropic. Measuring AI agent autonomy in practice. 18 Feb 2026.
\bibitem{fermat} Zhang et al. Safely scaling multi-agent teams with Claude Fable 5.2. claude.dev, 9 Nov 2026.
\bibitem{harness} P. Rajasekaran. Harness design for long-running application development. Anthropic, 24 Mar 2026.
\bibitem{longrun} J. Young. Effective harnesses for long-running agents. Anthropic, 26 Nov 2025.
\bibitem{zhou} Zhou et al. When Should Users Check? CHI 2026. arXiv:2510.05307
\bibitem{wynn} Wynn, Satija and Hadfield. Talk Isn't Always Cheap. arXiv:2509.05396
\bibitem{cais} Bertalani\v{c} and Fortuna. The Cost of Consensus. ACM CAIS 2026.
\bibitem{pappu} Pappu et al. Multi-Agent Teams Hold Experts Back. arXiv:2602.01011v4
\bibitem{carlini} N. Carlini. Building a C compiler with a team of parallel Claudes. Anthropic, 5 Feb 2026.
\bibitem{metrrisk} METR. Frontier Risk Report (Feb to Mar 2026). 19 May 2026.
\bibitem{cursor} W. Lin. Scaling long-running autonomous coding. Cursor, 14 Jan 2026.
\bibitem{mckinsey} Krivkovich et al. AI is changing work. Now it has to change the organization. McKinsey, 5 Oct 2026.
\bibitem{bcgmit} MIT Sloan Management Review and BCG. The Emerging Agentic Enterprise. 18 Nov 2025.
\bibitem{masters} Masters et al. Orchestrating Human-AI Teams: The Manager Agent. DAI 2025. arXiv:2510.02557
\bibitem{metrrh} METR. Recent Frontier Models Are Reward Hacking. 5 Jun 2025.
\bibitem{agentindex} Staufer et al. The 2025 AI Agent Index. arXiv:2602.17753
\end{thebibliography}}
\end{multicols}
\end{document}
